\begin{abstract}

This report addresses the structural optimisation of a natural-draught cooling tower, modelled as a vertical stack of $m = 10$ conical frustums. 
The objective is to minimise the total lateral surface area, a proxy for construction material cost, 
subject to a fixed internal volume constraint of $V_{\text{target}} = 70{,}320\,\text{m}^3$. 
The equality constraint is handled via a quadratic penalty method, enabling a uniform, solver-agnostic problem formulation.

Three optimisation algorithms are benchmarked: the quasi-Newton Broyden-Fletcher-Goldfarb-Shanno (BFGS) method, 
Particle Swarm Optimisation (PSO), and Simulated Annealing (SA). 
Their performance is evaluated across eight study cases of increasing complexity, ranging from a fixed-height reference case ($D = 9$) 
to a parametric hyperboloidal fit ($D = 3$), 
encompassing variations in dimensionality, initialisation strategy, constraint topology, and physical scaling.

Results demonstrate that NEED TO GET THE FINAL RESULTS 
\end{abstract}
